
This study quantified, at population scale across 223 diverse LLMs, the degree to which model capability independently predicts length-debiased preference after controlling for response verbosity. The results are unambiguous: a partial Spearman correlation of 0.985 (p = 1.69e-170, one-tailed, bootstrap 95\% CI [0.976, 0.988]) substantially exceeds the bivariate agreement reported in prior work. Two evaluation systems --- human pairwise comparison (win\_rate) and GLM-debiased annotation (LC\_winrate) --- agree on model capability ordering at a level that has not been previously quantified.

The principal contributions are: (1) r\_partial = 0.985 (p = 1.69e-170, one-tailed, N = 223) establishing population-scale capability-LC alignment after verbosity control; (2) capability-verbosity dominance 4.88:1 in standardized OLS, with a negative verbosity coefficient ($\beta _len = -4.37)$ confirming LC correction functionality and explaining 70.7 additional percentage points of $R^2$ beyond verbosity alone; (3) FWL-inspired verification within 1.1 percentage points ($\rho_{resid}$ = 0.9739), ruling out methodological artifact; (4) monotonic quartile ordering with large effect ($\epsilon ^2$ = 0.883, Dunn Q1 vs Q4 p = 1.04e-38); and (5) a $10\times$ dependent variable choice effect-size lesson favoring LC\_winrate over $\Delta$.

Future work should pursue: cross-framework replication (MT-Bench, Chatbot Arena, N $\geq$ 100 model overlap); family-stratified partial correlation analysis to assess whether the finding is driven by model family clustering; longitudinal extension to post-2024 frontier models; Pearson-based FWL exact equality verification; and investigation of why mid-capability (Q3) models exhibit the largest median alignment gap.

The LC correction preserves capability ordering as operationalized by human preference win\_rate in the AlpacaEval 2.0 leaderboard. Future evaluation research can use this result as an empirical baseline when choosing between raw and length-controlled preference metrics.

